% GENERATED by scripts/build.py convert ch01 -- do not hand-edit; edit
% chapters/md/01_introduction.md and regenerate.
% md sha256 b121a129a694669260bee66d3cb8e9ff6d4636434ebde306b00d725a82677862
% canonical manuscript/v13_2026-10-17/01_introduction.md sha256 b121a129a694669260bee66d3cb8e9ff6d4636434ebde306b00d725a82677862
% md identical to canonical: True
% filter sha256 5f7f0d210d735950ef97c22d7ae51d76d95fce54efc094b16108921e53a47b1c

\section{Introduction}\label{sec:1}\ThesisCheckNumber{section}{1}

Child stunting remains one of the most persistent markers of population health disadvantage in Sub-Saharan Africa. Linear growth faltering is concentrated in the first two years of life, is difficult to reverse once established, and its consequences extend well beyond anthropometry to cognitive development, school attainment, and adult earning capacity \autocite{Victora2008,Victora2010,Black2013}. Among the environmental exposures proposed to explain why children in similar economic circumstances grow at such different rates, water, sanitation, and hygiene (WASH) conditions occupy a central place in both the public-health and the development-economics literature, via several biologically plausible mechanisms supported to varying degrees: inadequate water and sanitation can raise exposure to faecal pathogens, increasing diarrhoeal disease, soil-transmitted helminth infection, and sub-clinical environmental enteric dysfunction (EED), each of which could divert nutritional intake toward immune response rather than growth \autocite{CummingCairncross2016} --- distinct, separately evidenced pathways, not one undisputed mechanism, with EED specifically flagged as the least settled and needing more research. Even the strongest recent causal evidence against basic WASH interventions reducing stunting does not dispute this plausibility (Section 2.3), but plausibility is not the same claim as an established causal mechanism, and this thesis does not conflate the two. What is considerably less settled is whether the mechanism translates into a \emph{detectable, policy-relevant statistical association} in observational survey data once household wealth and local development are accounted for, and whether a \emph{basic facility-level} WASH classification, as opposed to a comprehensive intervention package, is the right exposure in which to look for it --- the narrower, empirical question this thesis addresses.

\subsection{Two distinct exposures: a household's own facilities and its surroundings}\label{sec:1.1}\ThesisCheckNumber{subsection}{1.1}

A household's own water source and sanitation facility are not the only WASH-related exposure a child faces. Even a child in a household with an improved latrine and piped water can be exposed to pathogens through shared water sources, open defecation by neighbours, flooding, or simple proximity to inadequately managed waste if the surrounding community lacks similar facilities. The distinction between \emph{own-household WASH}, whose protective benefit accrues most directly to the household's own members, and \emph{surrounding-community WASH}, a local environmental condition shaped by the facilities of all households together, corresponds to two different exposure pathways --- it is not a modelling convenience, nor a claim that household WASH is a purely private good with no effect beyond its own household. Improved household sanitation in particular is a textbook source of a local environmental externality, since it reduces the pathogen load the household itself releases into the shared local environment. Recent community-level WASH research \autocite{Donohue2023,SkoufiasVinha2026} suggests that household and community WASH need not move together or carry the same association with child growth. For this reason, the thesis constructs and compares both measures: a child's own household classification, and a leave-one-out measure of improved-facility coverage among all \emph{other} households in the same DHS cluster, so that a child's own household is never used to construct its own community exposure.

\subsection{Exposure timing in cross-sectional survey designs}\label{sec:1.2}\ThesisCheckNumber{subsection}{1.2}

A specific measurement limitation affects the cross-sectional, repeated-survey-round observational literature to which the baseline analysis of this thesis belongs, including DHS-based studies \autocite{Rakotomanana2020} and the related national household-survey literature such as Addae et al.'s \autocite*{Addae2024} Ghana MICS-based study. In these designs, household and community WASH status is recorded once, at the survey date, and applied to every child in that round regardless of the child's age: a four-year-old measured in a 2022 survey experienced their nutritionally critical first two years of life around 2018--2019, under conditions the survey-date measurement cannot describe, whereas an infant measured in the same round has WASH recorded very close to their own early life. This does not apply to the major WASH-intervention trials, which assign exposure near each child's own birth and follow growth prospectively over a defined window (Section 2.3), nor to panel designs that track indicators over successive rounds within the same geographic units (Section 2.2). Age controls and age-interaction tests, also estimated here, describe how the \emph{current} association varies with how old a child is today; they do not reconstruct what a particular child's own environment actually was during their developmental window. The baseline analysis documents this limitation explicitly and does not claim to resolve it.

A separate, clearly bounded extension responds to one specific dimension of this limitation. It does not ask whether assigning a \emph{historical} instead of a \emph{current} exposure changes an otherwise identical estimate, because the baseline and historical samples differ too much in other respects to isolate that comparison (Section 6.2). Instead, it asks a self-contained question within its own sample: \emph{is modelled local WASH coverage during a child's own early-life window associated with that child's later anthropometric outcomes?} It matches each child's birth-cohort location to a multi-year, modelled local WASH coverage surface for a defined early-life window, using IHME's gridded 2000--2017 drinking-water and sanitation access estimates \autocite{LBDWaSH2020}, providing genuine within-cluster, cross-birth-cohort variation that a single survey round cannot offer. This improves measurement on one specific dimension --- an exposure closer to each child's developmental period rather than to the survey date --- without resolving the exposure-timing problem in general or reconstructing any household's actual historical WASH facilities, which no available data source can provide. The extension has its own separate limitations (undocumented measurement-subsample weighting, DHS/IHME input overlap, residence as a proxy for historical residence, confidentiality-driven coordinate displacement), stated in full in Section 6, which govern how much weight its results can bear.

\subsection{Research questions}\label{sec:1.3}\ThesisCheckNumber{subsection}{1.3}

This thesis asks three related questions, each posed and answered as an \emph{association under a cross-sectional design}, not as a causal effect --- stated once here, governing how every coefficient in this thesis should be read. First, is a child's own household's water and sanitation status associated with their anthropometric outcomes (HAZ, WAZ, WHZ) after accounting for household- and cluster-level confounding? Second, is the WASH status of the surrounding community, independent of the child's own household, separately associated with these outcomes? Third, within the three countries of the historical sample, is modelled local WASH coverage during a child's specific early-life window associated with these outcomes? The third question is answered only within the historical sample; its answer is not compared with the first two as if the comparison isolated the effect of exposure timing (Section 1.2; Section 6.2).

\subsection{Scope: four countries, two time horizons}\label{sec:1.4}\ThesisCheckNumber{subsection}{1.4}

The baseline analysis pools the most recent Demographic and Health Survey (DHS) round available, as of the data-collection date, for four countries: Ethiopia (2024--25), Ghana (2022), Kenya (2022), and Nigeria (2024), selected for data availability and recency, not as a representative sample of Sub-Saharan Africa; no claim in this thesis generalises findings from these four countries to the region as a whole. The historical extension uses an earlier DHS round for three of the four countries (Ghana 2014, Kenya 2014, and Nigeria 2018), linked to IHME's annual gridded WASH-coverage estimates, available through 2017. Ethiopia's 2016 round is retained only as a separately labelled, provisional sensitivity analysis, because the conversion of Ethiopian-calendar interview dates, which governs eligibility for each exposure window, could not be fully validated (Section 3.5 and Appendix C). The baseline and historical samples are \emph{not the same children}. They use different exposure definitions and different spatial and temporal constructions, and they are analysed, validated, and reported separately throughout. No result in one sample is used to adjust, select, or validate a result in the other.

\subsection{Contribution}\label{sec:1.5}\ThesisCheckNumber{subsection}{1.5}

This thesis makes two separate contributions in two separate samples, and they should not be read as one combined design. In the \emph{baseline sample}, the contribution is the explicit, side-by-side construction and comparison of own-household and leave-one-out community WASH exposure within one four-country sample and one consistent set of controls, with each exposure estimated under the fixed-effects structure its own construction requires (Sections 4.2--4.3). In the \emph{historical sample}, the contribution is a transparently documented, preliminary attempt to introduce genuine temporal variation in early-life exposure through geospatially matched historical coverage data. This extension examines a single modelled local-coverage exposure assigned to each child's own early-life window (Sections 3.5--3.8). It does not reproduce the household-versus-community comparison: the historical data source does not distinguish a child's own household from its surrounding community, so there is no household/community contrast to draw within the historical sample. A particular combination of four countries is not, by itself, evidence of substantive novelty, and this thesis does not claim it as such. Each of the design choices above is reported with the methodological limitations of its own sample, and these limitations are treated as integral to its interpretation rather than as closing caveats.

\subsection{A restrained summary of findings}\label{sec:1.6}\ThesisCheckNumber{subsection}{1.6}

The main findings are summarised briefly here; Sections 5--6 develop them with full uncertainty, and the Abstract gives the exact coefficients. Household-level WASH shows no statistically significant association with any outcome once cluster fixed effects are used to compare households within the same community. Community-level WASH coverage shows a large, strongly significant unadjusted association for most outcome/exposure pairs, which attenuates to statistical insignificance once socioeconomic controls and region fixed effects are added. This pattern is consistent with, but not conclusive proof of, confounding by local socioeconomic conditions. Within the historical sample, the pooled established-date water--HAZ association is positive and significant, but sensitive to the pooled weighting convention and carries a substantial list of unresolved documentation issues (Section 6; Appendix C). These results do not identify a causal policy effect or cost-effectiveness; they describe conditional statistical associations in specific survey samples, which can inform research questions and interpretation without themselves supporting a policy recommendation.

\subsection{Design and principal limitations, stated at the outset}\label{sec:1.7}\ThesisCheckNumber{subsection}{1.7}

This is a cross-sectional, associational design throughout. It is not a panel, difference-in-differences, or instrumental-variables design, and none of those terms is used elsewhere in this thesis to describe a method that is not actually implemented. The principal limitations are unresolved sampling-weight documentation for the measured-anthropometry subsample of the historical extension, DHS/IHME input overlap in two of the three historical-extension countries, the use of current residence as a proxy for historical residence, confidentiality-driven displacement of DHS cluster coordinates, and the small number of countries in both samples. Section 6 states them in full and does not treat them as a closing formality.

\subsection{Structure}\label{sec:1.8}\ThesisCheckNumber{subsection}{1.8}

Section 2 situates this thesis in the literature, distinguishing biological plausibility from whether a facility-level classification shows a detectable association, and causal evidence from conditional association. Section 3 describes the data and variable construction for both samples; Section 4 the empirical strategy; Section 5 the results. Section 6 relates the findings to the literature, explains why the baseline and historical results are not comparable as a test of exposure timing, and states all limitations in full. Section 7 concludes.
